\begin{tikzpicture}[
    >=Latex,
    node distance=0pt,
    every node/.style={transform shape},
    stage/.style={
      draw=villatheme, line width=0.9pt, fill=villatheme!6,
      rounded corners=3pt, align=center, inner sep=6pt,
      minimum width=2.7cm, minimum height=1.3cm, font=\small
    },
    flow/.style={->, line width=1pt, draw=villatheme},
    back/.style={->, line width=1pt, draw=villaaccent, dashed}
  ]

  % 四个环节，顺时针：上→右→下→左
  \node[stage] (write)   at ( 0   , 2.3) {\textbf{写 .tex}\\[3pt]{\scriptsize 这一页的内容}};
  \node[stage] (compile) at ( 4.6 , 0  ) {\textbf{编译}\\[3pt]{\scriptsize xelatex}};
  \node[stage] (render)  at ( 0   ,-2.3) {\textbf{渲染 PNG}\\[3pt]{\scriptsize pdftoppm}};
  \node[stage] (look)    at (-4.6 , 0  ) {\textbf{亲眼看}\\[3pt]{\scriptsize Read 渲染结果}};

  % 顺时针主回路（弧线相连）
  \draw[flow] (write)   to[bend left=18] (compile);
  \draw[flow] (compile) to[bend left=18] (render);
  \draw[flow] (render)  to[bend left=18] (look);
  \draw[flow] (look)    to[bend left=18] (write);

  % 不满意 → 改：从「看」回到「写」的闭环回箭头
  \draw[back] (look) to[bend right=42] (write)
    node[midway, font=\scriptsize, text=villaaccent,
         fill=white, inner sep=2pt] {不满意 → 改};

\end{tikzpicture}
